%%% CHANGELOG
%%% Patch 0605a (Session 146, 2026-05-27)
%%% Umbrella registration of THEO-DSL-7: Edge-Aligned O(delta^2) Closure Theorem.
%%%
%%% Combines into a single citable theorem package:
%%%   - Patch 0600 (edge-aligned first-order structural, hypothesis H1_E)
%%%   - Patch 0605 (edge-aligned second-order coefficient closure, hypothesis H5_E)
%%%
%%% Rigor level: sketch-document Layer 3 (matches constituents).
%%% Single-artifact discipline: 21-for-21 since Patch 0585.

\documentclass[11pt]{article}
\usepackage{amsmath,amssymb,amsthm}
\usepackage{geometry}
\geometry{margin=1in}

\newtheorem{theorem}{Theorem}
\newtheorem{corollary}{Corollary}
\newtheorem{remark}{Remark}

\title{THEO-DSL-7: Edge-Aligned $\mathcal{O}(\delta^2)$ Closure Theorem \\[2pt]
\large Umbrella Registration (Patches 0600 + 0605)}
\author{Hyperphysics Institute --- Conscious Point Physics Programme}
\date{Patch 0605a \quad Session 146 \quad 2026-05-27}

\begin{document}
\maketitle

\begin{abstract}
This document registers THEO-DSL-7 as the umbrella theorem covering the
edge-aligned current expansion through second order in $\delta$. The
combined statement, drawn from Patch 0600 (first order) and Patch 0605
(second order), asserts that at edge-aligned orientation $\theta_E$, the
$D_5$-invariant current $\mathbf{j}(\theta_E, \delta)$ is confined
order-by-order to the 2D non-orthogonal subspace
$\mathrm{span}\{\hat{\mathbf{n}}_\rho, \hat{\mathbf{n}}_{\text{edge}}\}$
through $\mathcal{O}(\delta^2)$, with explicit, real, finite,
non-vanishing coefficients at both orders. Four umbrella-level findings
(F-DSL-7-1 through F-DSL-7-4) record cross-patch observations not visible
at the level of either constituent alone.
\end{abstract}

\section{Theorem Statement (THEO-DSL-7)}

\begin{theorem}[Edge-Aligned $\mathcal{O}(\delta^2)$ Closure --- THEO-DSL-7]
\label{thm:THEO-DSL-7}
Let $\mathbf{j}(\theta_E, \delta)$ denote the $D_5$-invariant current at
edge-aligned orientation, expanded as
\[
\mathbf{j}(\theta_E, \delta) \;=\; \mathbf{j}_1\,\delta \;+\; \mathbf{j}_2\,\delta^2
\;+\; \mathcal{O}(\delta^3).
\]
Then through $\mathcal{O}(\delta^2)$:
\begin{enumerate}
\item[\textnormal{(i)}] \textbf{2D confinement (hypotheses H1\_E and H5\_E).}
Both $\mathbf{j}_1$ and $\mathbf{j}_2$ lie within
$\mathrm{span}\{\hat{\mathbf{n}}_\rho, \hat{\mathbf{n}}_{\text{edge}}\}$
and admit non-orthogonal-basis decompositions
\begin{align}
\mathbf{j}_1 &\;=\; \alpha_{1,\rho}\,\hat{\mathbf{n}}_\rho
              \;+\; \alpha_{1,\text{edge}}\,\hat{\mathbf{n}}_{\text{edge}}, \\
\mathbf{j}_2 &\;=\; \alpha_{2,\rho}\,\hat{\mathbf{n}}_\rho
              \;+\; \alpha_{2,\text{edge}}\,\hat{\mathbf{n}}_{\text{edge}}.
\end{align}

\item[\textnormal{(ii)}] \textbf{First-order coefficients} (Patch 0600
Theorem 4.1): $\alpha_{1,\rho}$ and $\alpha_{1,\text{edge}}$ are real,
finite, and non-vanishing (Finding A1-F1).

\item[\textnormal{(iii)}] \textbf{Second-order coefficients} (Patch 0605
Theorem 1):
\[
\alpha_{2,\rho} \;=\; \frac{9}{2\varphi} \;=\; \frac{9(\varphi-1)}{2}
\;\approx\; +2.7812,
\qquad
\alpha_{2,\text{edge}} \;=\; 9\varphi - 12 \;=\; 3(3\varphi-4)
\;\approx\; +2.5623.
\]
Both are positive (contrasting vertex-aligned analog $\alpha_2 = -9/\varphi^2 < 0$).

\item[\textnormal{(iv)}] \textbf{Second-order squared magnitude} (Patch 0605):
\[
|\mathbf{j}_2|^2 \;=\; \frac{9(184 - 111\varphi)}{4}\,(r_0\delta^2)^2
\;\approx\; 9.896\,(r_0\delta^2)^2.
\]
\end{enumerate}
\end{theorem}

\medskip
The rigor of THEO-DSL-7 inherits from its constituents: sketch-document
Layer 3 under (H1\_E) and (H5\_E), per DG-F15A-5 default-(A). No new
proof obligations are introduced by the umbrella; this registration
provides a single citable theorem statement for external reference.

\section{Constituent Patches and Results}

\subsection*{Patch 0600 --- Edge-Aligned First-Order Structural}
\begin{itemize}
\item Hypothesis (H1\_E): $\mathbf{j}_1 \in
\mathrm{span}\{\hat{\mathbf{n}}_\rho, \hat{\mathbf{n}}_{\text{edge}}\}$.
\item Gram matrix of the non-orthogonal basis established (§3.1).
\item Coefficients $\alpha_{1,\rho}, \alpha_{1,\text{edge}}$ derived (§4).
\item Finding A1-F1: 2D subspace genuinely 2-dimensional at
$\mathcal{O}(\delta)$.
\item File: \texttt{hardened\_theorems/edge\_aligned\_invariant\_subspace\_structural.tex}
\end{itemize}

\subsection*{Patch 0605 --- Edge-Aligned Second-Order Closure}
\begin{itemize}
\item Hypothesis (H5\_E): $\mathbf{j}_2 \in
\mathrm{span}\{\hat{\mathbf{n}}_\rho, \hat{\mathbf{n}}_{\text{edge}}\}$.
\item Inline Lemma: cross-shell $(S_1 \to S_3)$ edge formula with
additive constant $+\varphi^2/2$.
\item Per-orbit (B.1)--(B.4) sub-class decomposition table.
\item Orbit-level $\Omega(O_i)$ synthesis, including exact
$\Omega(O_2) = 0$.
\item Coefficients $\alpha_{2,\rho}, \alpha_{2,\text{edge}}$ closed.
\item Findings A5-F1 through A5-F5.
\item File: \texttt{hardened\_theorems/edge\_aligned\_second\_order\_alpha2.tex}
\end{itemize}

\section{Umbrella-Level Findings}

These findings record cross-patch observations not isolated at the level
of either constituent alone.

\paragraph{F-DSL-7-1: Order-by-Order 2D Confinement.}
The same 2D $D_5$-invariant subspace
$\mathrm{span}\{\hat{\mathbf{n}}_\rho, \hat{\mathbf{n}}_{\text{edge}}\}$
confines the current at \emph{both} $\mathcal{O}(\delta)$ and
$\mathcal{O}(\delta^2)$. A priori, second-order corrections could escape
the first-order invariant subspace via the symmetric tensor square; that
they do not is a non-trivial consequence of the $D_5$ representation
content of the path-class structure.

\paragraph{F-DSL-7-2: Path-Class Weights are Orientation-Independent.}
The path-class weights
$\{p_1, p_2, p_3, p_4\} = \{1,\, 1/2,\, -1/(2\varphi),\, -\varphi/2\}$
established in Patch 0591 §3 transfer unchanged from vertex-aligned to
edge-aligned orientation. What varies between orientations is the
projection functional $S_{\text{edge}}(O_i)$, not the weights.

\paragraph{F-DSL-7-3: B.2 Sub-Class as Orientation Discriminant.}
The intra-shell cross-vertex sub-class $B_2(O_i)$ vanishes identically at
vertex-aligned orientation (Patch 0591 §4.3) but is non-vanishing at
edge-aligned orientation (Patch 0605 Eq.~5, Remark 4.1). This is the
structural discriminant between the two orientations and the proximate
cause of the sign reversal in $\alpha_2$ (Finding A5-F2).

\paragraph{F-DSL-7-4: Edge-Aligned $\mathcal{O}(\delta^2)$ Program Complete.}
With THEO-DSL-7, the edge-aligned current is fully determined through
second order in $\delta$ within the path-class framework. The
next-order extension to $\mathcal{O}(\delta^3)$ is registered as open
(OPEN-A5-2 in Patch 0605; OPEN-DSL-7-2 below).

\section{Registration Record}

\begin{center}
\begin{tabular}{ll}
\hline
\textbf{Registry identifier:} & THEO-DSL-7 \\
\textbf{Series:} & THEO-DSL (theorem-level package registrations) \\
\textbf{Constituent patches:} & 0600, 0605 \\
\textbf{Hypotheses covered:} & (H1\_E), (H5\_E) \\
\textbf{Rigor level:} & sketch-document Layer 3 (matches constituents) \\
\textbf{Governance designation:} & DG-F15A-5 default-(A) \\
\textbf{Companion (vertex-aligned $\mathcal{O}(\delta^2)$):} & none yet
\\
\textbf{Registration session:} & 146 \\
\textbf{Registration date:} & 2026-05-27 \\
\textbf{Single-artifact discipline:} & 21-for-21 since Patch 0585 \\
\hline
\end{tabular}
\end{center}

\section{Citation Form}

When citing this package from external papers (CPP series and beyond),
use the form:

\medskip
\noindent\texttt{[THEO-DSL-7]} \emph{Edge-Aligned $\mathcal{O}(\delta^2)$
Closure Theorem.} CPP Programme internal registry, Session 146,
2026-05-27. Constituents: Patch 0600 (first-order structural) and
Patch 0605 (second-order coefficient closure).

\section{Forward Items}

\begin{itemize}
\item OPEN-DSL-7-1: Vertex-aligned $\mathcal{O}(\delta^2)$ umbrella
analog. Would combine Patch 0591 (vertex-aligned second-order) with a
not-yet-extant first-order vertex-aligned structural patch. Candidate
identifier: THEO-DSL-8.
\item OPEN-DSL-7-2: $\mathcal{O}(\delta^3)$ extension at edge-aligned
orientation. Currently OPEN-A5-2 in Patch 0605.
\item OPEN-DSL-7-3: Cross-orientation interpolation theorem
parameterizing $\alpha_2(\theta)$ between $\theta_V$ and $\theta_E$,
exploiting F-DSL-7-3 (B.2 sub-class as orientation discriminant).
\end{itemize}

\end{document}
